\documentclass[11pt,a4paper]{article}
\usepackage[utf8]{inputenc}
\usepackage[T1]{fontenc}
\usepackage{amsmath,amssymb}
\usepackage{booktabs}
\usepackage{geometry}
\usepackage{hyperref}
\usepackage{siunitx}
\usepackage{pifont}
\geometry{margin=2.5cm}

\title{Modellazione delle prestazioni di un DBMS\\ a partire da una traccia di workload}
\author{Marco Chiodo, Giuseppe Casteltrione}
\date{Relazione di progetto 3 -- Tecnologie e Architetture per la Gestione dei Dati}

\begin{document}
\maketitle

\section{Introduzione}

La progettazione e il dimensionamento dell'infrastruttura hardware di un DBMS richiedono
l'utilizzo di strumenti quantitativi in grado di stimare, già in fase preliminare, l'impatto che
eventuali variazioni del carico di lavoro o della configurazione del sistema possono avere sulle
prestazioni osservate dagli utenti.

La teoria delle code rappresenta un approccio consolidato per l'analisi di tali problematiche,
poiché consente di modellare il comportamento del sistema a partire da parametri misurabili,
come l'utilizzo delle risorse, i tempi di servizio e le caratteristiche del traffico generato. Attraverso
la costruzione di un modello analitico è quindi possibile: valutare le prestazioni del sistema, confrontare
le previsioni con i dati rilevati e supportare le decisioni di capacity planning, riducendo la necessità
di effettuare modifiche dirette sull'ambiente di produzione con rischi e costi elevati.

In questo lavoro si applica tale approccio alla modellazione delle prestazioni di un DBMS, a partire
da una traccia di workload reale (\texttt{transactions.csv}) raccolta osservando il sistema per un
periodo di 24 ore, durante il quale sono state registrate $N=100\,000$ transazioni. Per ciascuna
transazione la traccia riporta il tempo di CPU richiesto, il numero di accessi a un array di dischi
in configurazione RAID-0 e la percentuale di tali accessi costituita da operazioni di lettura.
Inoltre, nella vengono rese note le utilizzazioni medie delle due risorse osservate nello
stesso intervallo di tempo: $U_{CPU}=23\%$ per la CPU e $U_{RAID}=19\%$ per l'array dischi. Riassumendo, il
progetto richiede di costruire un modello di teoria delle code del DBMS, calibrato sulla traccia di
workload reale, per ricavare le prestazioni attuali del sistema, il carico massimo sostenibile e
l'impatto di un cambio di configurazione dell'array dischi da RAID-0 a RAID-1.

\section{Numero di core della CPU}

Il corretto dimensionamento del numero di core del centro CPU è un passaggio preliminare
imprescindibile: il numero di serventi non è solo un parametro architetturale, ma condiziona
direttamente come vanno interpretate le utilizzazioni fornite dal progetto e come va calibrata la
service demand nelle fasi successive.

\subsection{Vincolo necessario: capacità elaborativa minima:}

Sommando i contributi individuali del tempo di CPU riportato nella traccia per tutte le
$100\,000$ transazioni si ottiene il carico complessivo di lavoro da smaltire:
\begin{equation}
B_{cpu} = \sum_i \text{CPUTime}_i \approx 180\,583\ \text{s} \approx 50.16\ \text{ore}
\end{equation}
Poiché un singolo core può erogare al più $T=24$ ore di servizio effettivo in una giornata, e il
carico osservato richiede oltre il doppio di tale capacità, è necessario un livello di parallelismo
maggiore. Un primo limite inferiore al numero di core, puramente di fattibilità fisica, è
quindi:
\begin{equation}
m_{\min} = \left\lceil \frac{B_{cpu}}{T} \right\rceil = \left\lceil \frac{50.16}{24} \right\rceil = 3
\end{equation}
Questo valore, tuttavia, garantisce solo che il sistema sia fisicamente realizzabile. Di conseguenza, non è
detto che $m=3$ sia il numero di core effettivamente installato, né che sia coerente con l'utilizzazione media del 23\% dichiarata per il progetto. Determinare $m$ unicamente da questo vincolo, senza un ulteriore confronto con i dati osservati, porterebbe a un dimensionamento privo di una validazione rassicurante.

\subsection{Determinazione del numero di core tramite calibrazione incrociata:}
\label{sec:core-calib}

Il numero di core effettivo può essere ricavato sfruttando due stime indipendenti della service demand $D_{cpu}$ e imponendo che siano coerenti tra loro.

La prima stima è puramente \emph{misurata}: il tempo medio di CPU per transazione riportato
nella traccia di workload:
\begin{equation}
\bar S_{cpu} = \frac{1}{N}\sum_i \text{CPUTime}_i \approx 1.8058\ \text{s/job},
\end{equation}
non dipende in alcun modo dal numero di core, è un dato di servizio grezzo, non filtrato dal
modello.

La seconda stima deriva invece dalla \textit{legge dell'utilizzazione per centri multi-servitore}. Il
throughput di sistema, calcolato dalla traccia come
\begin{equation}
X_0 = \frac{N}{T} = \frac{100\,000}{86\,400} = 1.157407\ \text{tx/s},
\end{equation}
e l'utilizzazione per-core $U_{CPU}=0.23$ fornita dal progetto sono legati alla service demand
dalla relazione
\begin{equation}
U_{CPU} = \frac{X_0 \cdot D_{cpu}}{m} \quad\Longrightarrow\quad D_{cpu}(m) = \frac{U_{CPU}\cdot m}{X_0}
\end{equation}
che dipende esplicitamente da $m$.

Le due grandezze, $\bar S_{cpu}$ e $D_{cpu}(m)$, rappresentano la stessa quantità fisica (la domanda
di servizio media della CPU per transazione) stimata con due metodi indipendenti: l'una dalla
traccia grezza di workload, l'altra dall'utilizzazione dichiarata. Il valore corretto di $m$ è quello che rende le
due stime coerenti:
\begin{equation}
\frac{U_{CPU}\cdot m}{X_0} = \bar S_{cpu} \quad\Longrightarrow\quad
m = \frac{\bar S_{cpu}\cdot X_0}{U_{CPU}} = \frac{1.8058 \times 1.157407}{0.23} \approx 9.087
\end{equation}
Arrotondando al numero intero più vicino si ottiene $m=9$, che rispetta ampiamente il vincolo di
fattibilità $m \ge m_{\min}=3$ individuato al passo precedente. Con $m=9$ la service demand
implicita risulta $D_{cpu}(9) = 0.23\times 9/1.157407 = 1.7885$ s/job, in accordo (a meno
dell'arrotondamento di $m$ a intero) con il valore misurato $\bar S_{cpu}=1.8058$ s/job.

È importante osservare che il criterio basato esclusivamente sul carico complessivo \\($m_{\min}=3$)
fornisce una condizione \emph{necessaria ma non sufficiente} per la fattibilità del sistema. Tale
valore rappresenta infatti solamente il numero minimo teorico di core richiesto per completare
l'elaborazione entro la finestra temporale considerata, ma non descrive la reale configurazione del
sistema. L'adozione di soli tre core porterebbe infatti a una service demand della CPU pari a
$D_{cpu}=0.596$ s/job, incompatibile con il tempo medio di CPU osservato nella traccia
($1.806$ s/job), con una discrepanza di circa un fattore $3$. Per tale motivo, nelle analisi
successive viene adottato il valore $m=9$ come numero effettivo di core del centro CPU, in quanto
coerente con i dati osservati e con l'utilizzazione dichiarata del sistema.

\section{Modello del sistema:}

Il sistema DBMS viene modellato mediante una rete di code aperta, composta da due centri di
servizio disposti in serie, corrispondenti alle due risorse hardware osservate:
\begin{itemize}
\item \textbf{CPU} -- centro multi-servitore costituito da $m=9$ core paralleli;
\item \textbf{Array RAID} -- centro multi-servitore che rappresenta il sottosistema di
memorizzazione, inizialmente configurato come RAID-0 a $n_{disk}=2$ dischi indipendenti.
\end{itemize}
Si assume che la stripe unit dell'array sia maggiore della dimensione di pagina del DBMS e che
gli accessi siano allineati, cosicché ogni visita logica all'array si traduca in un accesso a un solo
disco fisico: è questa ipotesi che rende sensato modellare l'array RAID-0 come una stazione a
$2$ serventi indipendenti anziché come un unico servente aggregato.

\subsection{Throughput di sistema:}

Il throughput di sistema è già stato ricavato dalla traccia nella Sez.~\ref{sec:core-calib}, in
fase di calibrazione del numero di core: $X_0 = N/T = 1.1574$ tx/s. Trattandosi di una misura indipendente dal numero di core, esso costituisce un parametro diretto del modello e viene
qui richiamato come base per il calcolo delle service demand nella sezione seguente.

\subsection{Service demands:}

Coerentemente con l'approccio adottato, le service demands non vengono ottenute come semplice
media dei tempi riportati nella traccia, ma tramite inversione della legge dell'utilizzazione,
usando il throughput e le utilizzazioni medie fornite dal progetto -- separando così il
dimensionamento hardware (numero di core, numero di dischi) dalla calibrazione del modello
prestazionale.

Per il centro CPU, con $m=9$ servitori paralleli:
\begin{equation}
D_{cpu} = \frac{U_{CPU}\cdot m}{X_0} = \frac{0.23 \times 9}{1.157407} = 1.7885\ \text{s/job}
\end{equation}
Per il centro RAID, con $n_{disk}=2$:
\begin{equation}
D_{raid} = \frac{U_{RAID}\cdot n_{disk}}{X_0} = \frac{0.19\times 2}{1.157407} = 0.3283\ \text{s/job}
\end{equation}
Dalla service demand del RAID è inoltre possibile ricavare il tempo di servizio del singolo
disco, dividendo per il numero medio di visite per transazione osservato in traccia
($\bar V_{raid}=2.7576$ accessi/transazione):
\begin{equation}
S_{raid} = \frac{D_{raid}}{\bar V_{raid}} = \frac{0.3283}{2.7576} = 0.1191\ \text{s}
\end{equation}
Questo valore di $S_{raid}$, calibrato sulla configurazione RAID-0, verrà riutilizzato al Punto 3 per
modellare il passaggio a RAID-1, poiché rappresenta una caratteristica fisica del singolo disco e
non della particolare organizzazione RAID adottata.

\section{Tempi di risposta e job nel sistema: }
In questa sezione vengono stimati il tempo medio di risposta del sistema e il numero medio di job presenti, due tra i principali indicatori delle prestazioni di un sistema a code. L'analisi viene condotta dapprima mediante un modello analitico semplificato, utilizzato come riferimento iniziale, e successivamente tramite una simulazione realizzata con JMT, in grado di rappresentare correttamente la natura multi-servente della CPU. Il confronto tra i due approcci consente di evidenziare i limiti dell'approssimazione analitica e di individuare una stima più accurata delle prestazioni effettive del sistema.
\subsection{Stima analitica di riferimento (limite superiore):}

Una prima stima, di solo riferimento, si ottiene applicando la formula per il tempo di risposta di
un centro a \emph{singolo} servente,
\begin{equation}
R_k = \frac{D_k}{1-U_k},
\end{equation}
sia alla CPU sia al RAID:
\begin{align}
R_{cpu} &= \frac{D_{cpu}}{1-U_{CPU}}
         = \frac{1.7885}{1-0.23}
         = 2.3227\ \text{s} \\
R_{raid} &= \frac{D_{raid}}{1-U_{RAID}}
          = \frac{0.3283}{1-0.19}
          = 0.4053\ \text{s} \\
R_{TOT} &= R_{cpu}+R_{raid}
         = 2.3227 + 0.4053
         = 2.7280\ \text{s} \\
N &= X_0\cdot R_{TOT}
  = 1.157407 \times 2.7280
  = 3.157\ \text{job}
\end{align}

Questa stima, però, sovrastima sistematicamente il tempo di risposta reale: la formula
$R=D/(1-U)$ presuppone che una richiesta si accodi non appena il servente è occupato, ipotesi
corretta per un centro a singolo servente ma non per un centro multi-servitore come la CPU con
$m=9$ core. Con 9 core al 23\% di utilizzazione, l'evento "tutti i 9 core occupati contemporaneamente"
è raro, per cui la formula a singolo servente è eccessivamente conservativa proprio sul centro
dominante del sistema.

\subsection{Validazione tramite simulazione JMT \textit{(JSIMgraph)}:}

Per ottenere una stima corretta si è simulato il modello -- rete aperta, Source $\to$ CPU (9
serventi) $\to$ RAID (2 serventi) $\to$ Sink, service demand $D_{i}$ per stazione -- con lo strumento
JSIMgraph della suite JMT. I risultati, confrontati con la stima analitica, sono riportati in
Tabella~\ref{tab:punto1}.

\begin{table}[h]
\centering
\caption{Confronto tra stima analitica (singolo servente) e simulazione JMT (multi-servente) -- Punto 1}
\label{tab:punto1}
\begin{tabular}{lrrr}
\toprule
 & Analitico $R=D/(1-U)$ & JMT (multi-server) & Scarto \\
\midrule
$X_0$ [tx/s] & 1.1574 & 1.1392 & 1.6\% \\
$R_{TOT}$ [s] & 2.7280 & 2.1235 & 28.5\% \\
$N$ [job] & 3.1574 & 2.4242 & 30.2\% \\
$U_{CPU}$ & 0.2300 & 0.2245 & 2.4\% \\
$U_{RAID}$ & 0.1900 & 0.1891 & 0.5\% \\
\bottomrule
\end{tabular}
\end{table}

Le utilizzazioni simulate (22.45\% CPU, 18.91\% RAID) riproducono, entro il rumore di
simulazione, quelle misurate dal performance monitor (23\%, 19\%): tali valori non sono stati
imposti ma emergono dalla parametrizzazione, il che valida il modello. Sul tempo di risposta,
invece, la formula a singolo servente e la simulazione divergono di circa il 28\%, per il motivo
discusso sopra.

Il dettaglio per componente conferma la spiegazione (Tabella~\ref{tab:punto1-percomp}): il tempo
di risposta della CPU simulato ($R_{cpu}=1.7829$ s) è quasi identico alla sua service demand
($D_{cpu}=1.7885$ s) -- la CPU, pur essendo il centro più caricato, non accoda praticamente nulla,
perché con 9 core al 23\% è estremamente improbabile che tutti i serventi siano occupati
contemporaneamente. Il RAID, pur meno utilizzato, ha solo 2 serventi e quindi accumula
un'attesa relativa maggiore rispetto alla propria service demand.

\begin{table}[h]
\centering
\caption{Tempi di risposta per componente -- Punto 1}
\label{tab:punto1-percomp}
\begin{tabular}{lrr}
\toprule
 & Analitico $R=D/(1-U)$ & JMT (multi-server) \\
\midrule
$R_{cpu}$ [s] & 2.3227 & 1.7829 \\
$R_{raid}$ [s] & 0.4053 & 0.3439 \\
Somma [s] & 2.7280 & 2.1268 \\
\bottomrule
\end{tabular}
\end{table}

La somma dei tempi per componente (1.7829+0.3439=2.1268 s) coincide, entro lo 0.16\%, con il
tempo di risposta di sistema misurato indipendentemente da JMT (2.1235 s), a conferma della
coerenza interna della simulazione.

\subsection{Risposta al Punto 1:}
Si assume come risultato del Punto 1 la stima JMT, corretta per la natura multi-servente della
CPU: tempo di risposta medio complessivo $R_{TOT}=2.12$ s ($R_{cpu}=1.78$ s, $R_{raid}=0.34$ s),
numero medio di job nel sistema $N=2.42$ (verificato con la legge di Little:
$N=X_0\cdot R = 1.1392\times 2.1235 = 2.42$).

\section{Carico massimo per-classe con $U_{bottleneck}\le 70\%$:}
Il secondo punto dell'analisi ha l'obiettivo di determinare il massimo carico di lavoro che il sistema può sostenere mantenendo l'utilizzazione del centro collo di bottiglia non superiore al 70\%. A tal fine, il workload viene dapprima suddiviso in classi omogenee mediante tecniche di clustering, così da verificare se le diverse tipologie di transazioni presentino domande di servizio differenti. Successivamente, utilizzando il modello prestazionale ricavato nei capitoli precedenti, viene calcolato il throughput massimo ammissibile e il corrispondente carico per ciascuna classe, verificando inoltre quale centro rimanga il collo di bottiglia del sistema al crescere del carico.
\subsection{Decomposizione in classi di workload:}
\label{sec:classi}

Analizzando la traccia con lo strumento di \emph{Workload Analysis} di JMT (clustering k-means a
2 cluster, $50\,000$ osservazioni ciascuno) emergono due classi di transazioni, di peso identico
$p_1=p_2=0.5$, ma con caratteristiche di accesso ai dischi molto diverse in termini di
letture/scritture:

\begin{table}[h]
\centering
\caption{Classi di workload individuate tramite clustering}
\begin{tabular}{lrr}
\toprule
 & Classe 1 & Classe 2 \\
\midrule
Peso $p_r$ & 0.50 & 0.50 \\
$V_r$ (visite I/O per transazione) & 2.7552 & 2.7599 \\
\%Read$_r$ & 10.02\% & 89.96\% \\
$D_{cpu,r}$ [s] & 1.7895 & 1.7874 \\
$D_{raid,r}$ [s] & 0.3280 & 0.3286 \\
\bottomrule
\end{tabular}
\end{table}

Le due classi hanno cioè domande di servizio praticamente identiche su CPU e RAID (differenza
$<0.2\%$): la classe 1 è a prevalenza di scritture (90\%), la classe 2 a prevalenza di letture (90\%).
Poiché i pesi sono 50/50 e le domande di servizio aggregate coincidono, il modello multi-classe
\emph{collassa} su quello mono-classe per tutto ciò che riguarda le utilizzazioni: un incremento
proporzionale del carico su tutte le classi non altera il mix, quindi $U_i$ scala dello stesso fattore
$\alpha$ indipendentemente dalla decomposizione in classi. La distinzione in classi diventerà invece
determinante al Punto 3, dove il rapporto letture/scritture, non l'utilizzazione aggregata,
condiziona il costo del mirroring.

Come verifica di coerenza, ricostruendo l'utilizzazione della CPU dal modello multi-classe si
ottiene $U_{CPU}=\sum_r D_{cpu,r}\cdot X_{0,r}/m = 0.2300$, identica al valore misurato ($X_{0,r}$ ottenuto come $X_{0,r} = X_0\cdot p_r$).

\subsection{Carico massimo sostenibile:}

Confrontando le utilizzazioni per-server ($U_{CPU}=0.23 > U_{RAID}=0.19$), la CPU resta il collo
di bottiglia del sistema. Uno scaling proporzionale del carico mantiene invariato il rapporto tra le
utilizzazioni dei due centri, per cui il fattore di scala massimo che porta la CPU al 70\% è:
\begin{equation}
\alpha = \frac{U_{\max}}{U_{CPU}} = \frac{0.70}{0.23} = 3.0435
\end{equation}
Equivalentemente, invertendo la legge dell'utilizzazione rispetto al throughput:
\begin{equation}
X_{0,\max} = \frac{U_{\max}\cdot m}{D_{cpu}} = \frac{0.70\times 9}{1.7885} = 3.5225\ \text{tx/s}
\quad (\approx 304\,348\ \text{transazioni/giorno})
\end{equation}
ripartito equamente tra le due classi: $X_{0,\max,r} = 1.7613$ tx/s ciascuna
($\approx 152\,174$ transazioni/giorno per classe).

A questo livello di carico l'utilizzazione della CPU raggiunge esattamente il 70\% (per
costruzione), mentre l'array RAID-0 sale a
\begin{equation}
U_{raid}(X_{0,\max}) = \alpha \times U_{RAID} = 3.0435\times 0.19 = 57.83\%,
\end{equation}
rimanendo ampiamente sotto soglia: la CPU resta quindi il collo di bottiglia anche al carico
massimo, confermando che la scelta di $U_{\max}$ come vincolo sulla CPU era corretta.

\section{Passaggio a RAID-1:}
Il terzo punto dell'analisi valuta l'impatto della sostituzione dell'attuale configurazione RAID-0 con un sistema RAID-1, mantenendo il throughput massimo determinato nel Punto 2 e imponendo che l'utilizzazione del centro disco non superi il limite del 70\%. A differenza dei punti precedenti, l'analisi deve tenere conto del diverso comportamento delle operazioni di lettura e scrittura in un array RAID-1: mentre le letture possono essere distribuite tra i dischi del mirror, le scritture devono essere replicate su ogni disco, modificando in modo sostanziale la service demand del centro. Sulla base di tale modello vengono quindi verificati la fattibilità della configurazione, il numero minimo di dischi necessario a rispettare il vincolo prestazionale e la coerenza dei risultati attraverso una simulazione realizzata con JMT.
\subsection{Modello del centro disco in RAID-1:}

Nel passaggio da RAID-0 (striping) a RAID-1 (mirroring a $n$ dischi) il comportamento del
centro disco cambia struttura:
\begin{itemize}
\item le \textbf{letture} possono essere servite da un disco qualunque del mirror: il carico di
lettura si distribuisce e si riduce al crescere di $n$;
\item le \textbf{scritture} devono essere replicate su tutti gli $n$ dischi del mirror: ogni disco
sostiene per intero il carico di scrittura, senza alcun beneficio di parallelismo.
\end{itemize}
Il tempo di servizio del singolo disco $S_{raid}=0.1191$ s, calibrato al Punto precedente sulla
configurazione RAID-0, resta valido: cambia il livello RAID (quante visite fisiche genera una
visita logica), non la tecnologia del disco. Senza questa ipotesi il Punto 3 non sarebbe risolvibile,
poiché per una configurazione non ancora esistente non si dispone di un'utilizzazione misurata da
cui ricavare i demand.

Separando dalla traccia le visite medie di lettura e scrittura per transazione,
\begin{align}
A &= \frac{1}{N}\sum_i \text{IO}_i\cdot\frac{\%\text{Read}_i}{100}
   = 1.3795\ \text{(letture/tx)} \\
B &= \frac{1}{N}\sum_i \text{IO}_i\cdot\left(1-\frac{\%\text{Read}_i}{100}\right)
   = 1.3781\ \text{(scritture/tx)}
\end{align}
con $A+B=2.7576$, coerente col numero medio di visite osservato. La ripartizione è quindi
pressoché paritaria ($\approx$50.03\%/49.97\%).

L'utilizzazione del centro RAID-1 a $n$ dischi, al throughput $X_0$, risulta:
\begin{equation}
U_{RAID\text{-}1}(n) = \frac{X_0\cdot S_{raid}\cdot(A + n\cdot B)}{n}
= \underbrace{\frac{X_0\cdot S_{raid}\cdot A}{n}}_{\text{letture}} + \underbrace{X_0\cdot S_{raid}\cdot B}_{\text{scritture}}
\label{eq:uraid1}
\end{equation}
Nel termine di scrittura la $n$ si semplifica: il lavoro si moltiplica per $n$ (replica su tutti i
dischi) e la capacità per $n$ (numero di serventi), per cui \textbf{aggiungere dischi non allevia mai
il carico di scrittura}, ma solo quello di lettura. Ne consegue un asintoto invalicabile
\begin{equation}
U_\infty = \lim_{n\to\infty} U_{RAID\text{-}1}(n) = X_0\cdot S_{raid}\cdot B
\end{equation}
la cui verifica ($U_\infty \le U_{\max}$) va effettuata \emph{prima} di cercare il valore di $n$, perché
se il solo carico di scrittura eccedesse la soglia nessun numero di dischi potrebbe soddisfare il
vincolo.

\subsection{Verifica di fattibilità e dischi necessari:}

Valutando il sistema al carico massimo trovato al Punto 2 ($X_{0,\max}=3.5225$ tx/s):
\begin{equation}
U_\infty = X_{0,\max}\cdot S_{raid}\cdot B = 3.5225\times 0.1191\times 1.3781 = 57.80\%
\end{equation}
Poiché $U_\infty = 57.80\% < 70\% = U_{\max}$, il vincolo di scrittura da solo \textbf{non preclude} la
fattibilità: esiste un numero finito di dischi che soddisfa il requisito. La feasibility non era
tuttavia scontata a priori: se la frazione di scritture del workload avesse superato il valore di
breakeven ($U_\infty=U_{\max}\Rightarrow$ frazione di scritture $\approx 60.5\%$ delle visite), nessun $n$
sarebbe stato sufficiente. Il workload analizzato ha il 49.97\% di scritture, sotto soglia ma non di
molto.

Applicando l'Eq.~\eqref{eq:uraid1} per $n=2,\dots,9$ dischi si ottiene la Tabella~\ref{tab:raid1}.

\begin{table}[h]
\centering
\caption{Utilizzazione del centro RAID-1 al variare del numero di dischi del mirror ($X_0=X_{0,\max}$)}
\label{tab:raid1}
\begin{tabular}{lrrrrrrrr}
\toprule
$n$ dischi & 2 & 3 & 4 & 5 & 6 & 7 & 8 & 9\\
\midrule
$U_{RAID\text{-}1}$ & 86.72\% & 77.08\% & 72.26\% & 69.37\% & 67.44\% & 66.06\% & 65.03\% & 64.23\%\\
Vincolo $\le 70\%$ & \ding{55} & \ding{55} & \ding{55} & \checkmark & \checkmark & \checkmark & \checkmark & \checkmark\\
\bottomrule
\end{tabular}
\end{table}

Il numero minimo esatto di dischi (a valori reali) si ottiene invertendo l'Eq.~\eqref{eq:uraid1}
rispetto a $n$:
\begin{equation}
n \ge \frac{X_0\cdot S_{raid}\cdot A}{U_{\max} - X_0\cdot S_{raid}\cdot B}
= \frac{3.5225\times 0.1191\times 1.3795}{0.70 - 0.5780} = 4.741 \quad\Longrightarrow\quad n_{\min}=\lceil 4.741\rceil = 5
\end{equation}
in accordo con la Tabella~\ref{tab:raid1}: $n=4$ non basta (72.26\%), $n=5$ soddisfa il vincolo
(69.37\%).

\subsection{Risposta al Punto 3:}
Sì, è possibile sostenere il throughput $X_{0,\max}$ con RAID-1 mantenendo l'utilizzazione
del collo di bottiglia $\le 70\%$: sono necessari \textbf{5 dischi} (mirror a 5 vie). La tolleranza ai
guasti richiesta (almeno 1 disco) è ampiamente soddisfatta, poiché un mirror a 5 vie tollera fino a
4 guasti simultanei.

Va tuttavia segnalato che la configurazione è \emph{al limite}: il margine sul vincolo è di soli 0.63
punti percentuali (69.37\% contro 70\%). Con 6 dischi il margine salirebbe a 2.56 punti percentuali
(67.44\%) al costo di un solo disco aggiuntivo -- una scelta più prudente in un contesto di
produzione reale.

Il costo della soluzione va inoltre osservato in termini di capacità utile: servono 5 dischi per un
array che, in RAID-0, ne richiedeva 2, con capacità utile pari a quella di un solo disco (fattore di
overhead di storage $5\times$). Il costo del mirroring, inoltre, non è distribuito uniformemente tra
le classi di workload: la classe 1 (90\% scritture) genera una service demand sul nuovo array
pari a $D_{RAID\text{-}1,1}=1.5087$ s, contro $D_{RAID\text{-}1,2}=0.4606$ s della classe 2 (90\%
letture) -- un rapporto di $3.28\times$, invisibile nei Punti 1 e 2 dove le classi erano
indistinguibili in termini di utilizzazione aggregata. È proprio la decomposizione in classi
effettuata al Punto 2 a rendere possibile attribuire questo costo differenziale.

\subsection{Validazione tramite simulazione JMT \textit{(JSIMgraph)}:}

Il modello RAID-1 a 5 dischi è stato inoltre simulato in JSIMgraph, con tasso di arrivo per classe
pari a $X_{0,\max,r}=1.7613$ tx/s e service demand per classe $D_{RAID\text{-}1,1}=1.5087$ s,
$D_{RAID\text{-}1,2}=0.4606$ s. Il confronto con il modello analitico è riportato in
Tabella~\ref{tab:punto3}.

\begin{table}[h]
\centering
\caption{Confronto tra modello analitico e simulazione JMT -- Punto 3 (RAID-1, 5 dischi)}
\label{tab:punto3}
\begin{tabular}{lrr}
\toprule
 & Analitico & JMT \\
\midrule
$X_0$ [tx/s] & 3.5225 & 3.5257 \\
$U_{CPU}$ & 0.7000 & 0.7063 \\
$U_{RAID}$ & 0.6937 & 0.7003 \\
\bottomrule
\end{tabular}
\end{table}

Gli scostamenti tra i due approcci (entro l'1\% su tutte le grandezze) sono imputabili al rumore
di simulazione e non a errori di parametrizzazione, come confermato ricalcolando i service
demand impliciti dai risultati JMT tramite la legge dell'utilizzazione: $D_{cpu}$ risulta 1.7736 s
(JMT) contro 1.7885 s (modello), $D_{raid}$ risulta 0.3320 s (JMT) contro 0.3283 s (modello).

Il passaggio a RAID-1 comporta inoltre un evidente degrado prestazionale rispetto al Punto 1: il
tempo di risposta di sistema passa da $R=2.12$ s a $R=3.21$ s ($+51\%$), e il numero medio di job
nel sistema da $N=2.42$ a $N=11.26$ -- conseguenza diretta dell'aver spinto entrambi i centri
vicino alla soglia del 70\% di utilizzazione.

\subsection{Osservazione conclusiva sul RAID-1:}

La radice del limite strutturale del RAID-1 è che le scritture si replicano su $n$ dischi e si
dividono su $n$ serventi: la $n$ si semplifica e resta un pavimento invalicabile
($U_\infty=57.80\%$ in questo caso). Una configurazione RAID-10 ($m'$ coppie mirrorate messe in
striping tra loro), replicando ogni scrittura su sole 2 unità e distribuendo le coppie via striping,
eliminerebbe questo asintoto e ridurrebbe ulteriormente l'utilizzazione per disco a parità di
dischi totali. Nel caso qui analizzato il vincolo del progetto impone però esplicitamente RAID-1,
e tale configurazione -- pur strutturalmente meno efficiente su un workload a scritture
significative come questo -- risulta comunque sufficiente con 5 dischi, senza necessità di ricorrere
al RAID-10.

\section{Implementazione in JMT:}

Il modello è stato realizzato utilizzando due moduli della suite JMT: il modulo di
\emph{Workload Analysis} per la decomposizione del carico in classi e il simulatore
\emph{JSIMgraph} per tutte le simulazioni (Punti 1 e 3).

\subsection{Workload Analysis -- decomposizione in classi:}
Il modulo di Workload Analysis è stato impiegato per la decomposizione del carico nelle due classi
già descritta nella Sez.~\ref{sec:classi}. Da quell'analisi provengono i parametri per-classe che
alimentano il simulatore -- i pesi $p_r$, il numero medio di visite $V_r$ e la frazione di letture
$\%\text{Read}_r$ -- dai quali sono state ricavate analiticamente le service demand per classe.

\subsection{JSIMgraph -- Punto 1 e Punto 3:}
Tutte le simulazioni adottano la stessa topologia, minimale: \emph{Source} $\to$ Queue ``CPU''
$\to$ Queue ``RAID'' $\to$ \emph{Sink}, con due classi aperte. Ogni stazione ha una sola visita e
tempo di servizio esponenziale pari direttamente alla service demand per classe $D_{i,r}$: ciò è
lecito perché $D_i = S_i V_i$ aggrega già tutte le visite nel demand, che diventa così l'input
diretto del modello (coerentemente con l'Esempio~8 delle slide).

\begin{itemize}
\item \textbf{Punto 1} (RAID-0): CPU con \emph{Number of Servers}$=9$, RAID con
\emph{Number of Servers}$=2$; arrival rate per classe $\lambda_r = X_0/2 = 0.5787$ tx/s; service
time $D_{cpu,r}=$ 1.7895 / 1.7874 s sulla CPU e $D_{raid,r}=$ 0.3280 / 0.3286 s sul RAID. La
simulazione restituisce $U_{CPU}\approx 0.22$ e $U_{RAID}\approx 0.19$: le utilizzazioni misurate
vengono riprodotte, non imposte, ed è questa la validazione del modello.
\item \textbf{Punto 3} (RAID-1, 5 dischi): stessa struttura, con la sola Queue ``RAID''
riconfigurata a \emph{Number of Servers}$=5$; arrival rate per classe
$\lambda_r = X_{0,\max}/2 = 1.7613$ tx/s; service time sul RAID pari alle service demand per classe
$D_{RAID\text{-}1,1}=1.5087$ s e $D_{RAID\text{-}1,2}=0.4606$ s. La CPU resta invariata.
\end{itemize}

Le due service demand del Punto 3 si ottengono analiticamente dai parametri della Workload
Analysis, tramite il modello del centro RAID-1 per classe con $n=5$ dischi:
\begin{equation}
D_{RAID\text{-}1,r} = S_{raid}\cdot V_r\cdot\big[\rho_r + n\,(1-\rho_r)\big],
\label{eq:draid1-perclasse}
\end{equation}
con $S_{raid}=0.1191$ s e $\rho_r$ frazione di letture della classe. Ne risulta
$D_{RAID\text{-}1,1}=0.1191\times 2.7552\times[0.1002+5(0.8998)]=1.5087$ s (classe write-heavy) e
$D_{RAID\text{-}1,2}=0.1191\times 2.7599\times[0.8996+5(0.1004)]=0.4606$ s (classe read-heavy).

La differenza di comportamento tra letture (servite da un solo disco) e scritture (replicate su
tutti gli $n$ dischi) \emph{non} è già
incorporata nei valori di service demand per classe tramite il fattore $[\rho_r + n(1-\rho_r)]$.
Fornendo al simulatore direttamente i $D_{RAID\text{-}1,r}$, la stazione multi-servitore a 5 dischi
riproduce correttamente l'utilizzazione del centro senza necessità di una topologia dedicata.

\section{Conclusioni:}

L'analisi svolta ha permesso di caratterizzare il comportamento prestazionale del sistema e di valutarne la capacità di sostenere incrementi di carico e modifiche architetturali. In primo luogo, la sola somma dei tempi di CPU presenti nella traccia, pari a circa 50.16 ore, dimostra che almeno tre core rappresentano un requisito necessario affinché l'elaborazione possa essere completata entro una finestra temporale di 24 ore. Tuttavia, imponendo la coerenza tra il tempo medio di CPU osservato nella traccia (1.806 s per transazione) e la service demand ricavata dall'utilizzazione dichiarata del 23\%, si ottiene una stima più realistica della configurazione del sistema, pari a nove core.

A partire dal throughput di sistema, pari a $X_0 = 1.1574$ transazioni al secondo, le service demand dei due centri sono stati determinati applicando la legge dell'utilizzazione, ottenendo $D_{cpu}=1.7885$ s e $D_{raid}=0.3283$ s. Tale approccio consente di ricavare la reale service demand, evitando di basarsi esclusivamente sui tempi medi grezzi riportati nella traccia CSV.

La modellazione del sistema tramite una rete di code, simulata mediante JMT, evidenzia un tempo medio di risposta pari a 2.12 s e un numero medio di job nel sistema di circa 2.42. Il confronto con il modello analitico a singolo servente mostra come quest'ultimo sovrastimi entrambe le grandezze di circa il 28--30\%, confermando che tale approssimazione non è adeguata per descrivere un centro multi-servitore poco carico come la CPU a nove core. Nonostante ciò, la CPU rimane il collo di bottiglia del sistema, con un'utilizzazione del 23\% per core, superiore a quella del RAID-0, pari al 19%.

L'analisi di scalabilità mostra inoltre che il sistema può sostenere un incremento uniforme del carico fino a circa 3.04 volte quello attuale, raggiungendo un throughput massimo di circa $3.52$ transazioni al secondo, corrispondente a circa 304,348 transazioni al giorno, senza superare il limite del 70\% di utilizzazione per core sulla CPU. In tali condizioni, il RAID-0 mantiene comunque un'utilizzazione inferiore alla soglia fissata, attestandosi al 57.83%.

Infine, è stata valutata la sostituzione del RAID-0 con una configurazione RAID-1. L'analisi dimostra che questa soluzione risulta sostenibile dal punto di vista prestazionale soltanto adottando un mirror composto da cinque dischi: una configurazione a due dischi porterebbe infatti l'utilizzazione fino all'86.7\%, superando il limite imposto. Con cinque dischi, invece, il sistema rispetta il vincolo del 70\% e garantisce la tolleranza al guasto di almeno un disco, arrivando a sopportarne fino a quattro. Tale soluzione comporta un costo di storage pari a cinque volte quello di un singolo disco e lascia un margine operativo piuttosto ridotto (0.63 punti percentuali), ma rende superfluo il ricorso a configurazioni RAID-10, inizialmente considerate in una valutazione più conservativa.

\end{document}